\section{Discussion}




We now discuss ongoing challenges, possible mitigations, and some deployment considerations, all staying within the general framework that we established for Passerine.









   





\subsection{Challenges and Possible Mitigations}\label{sec:challenges}

\begin{itemize}
    \item Context size: Due to verbose outputs (e.g. test logs, large source code files), and Passerine’s append-only history, the agent’s prompt to the LLM can exceed even large context windows (e.g., 2M tokens).
    We plan to explore possible mitigations including a SWE-Agent-style sliding window which elides the outputs of steps beyond the last n; the use of an LLM to summarize history; and sub-agents with their own, self-contained, history.

    
    \item Bug reproduction: Machine-reported bugs, like those produced by SAN or TOD, typically include some bug reproduction information. We observe that the agent typically uses this information to reproduce the bug as one of its first steps. In contrast, human-reported bugs include reproduction information at much lower rates. As a result, the agent must first figure out how to expose the bug. We have found that this is itself an interesting area of research, and one that would allow the agent to more effectively perform localization/editing. As a result, we have an ongoing project investigating the ability of an agent to generate bug-reproducing tests.
    
    \item Agent tools: 
    When we inspect agent responses that result in a tool usage failure, we find that for some bugs (particularly
    human bugs), the agent attempts to make use of non-existent commands. One prominent example is
    the agent attempting to use a \texttt{print} command to pose a question to the user
    (despite not having interactive capabilities). For example, in one run Passerine asks
    the user \emph{``print(Please provide more context about the issue with [...] What type of artifacts are missing?''}. Similarly, we observe Passerine can struggle to make progress
    on bugs that would benefit from 
    accessing documents/links listed
    in the bug description and the agent explicitly mentions 
    \emph{``i should try to access it''} (which we do not support).  We plan to extend Passerine's command set
    based on this observed behavior.

    \item Fault localization: We found that for human-reported bugs, fault localization can be challenging -- as reflected by low rates of successful file-level localization
    for trajectories that fail to produce a plausible patch. However, we believe
    there is substantial room for improvement. Currently, Passerine does not explicitly
    task the agent with localization and instead allows the dynamic loop to 
    perform steps as necessary.  In the future, iterative use of code search might be able to provide better localization, as has been argued in the AutoCodeRover~\cite{autocoderover} paper. Furthermore, our
    current implementation of Passerine does not exploit rich development histories
    which may reflect similar past bugs and patches, which can be used
    to further refine possible fault locations~.
    
    \item Diversity of patches: Passerine currently
    samples trajectories independently. Recent work~\cite{chatrepair} showed
    that conditioning LLM generation of patches on earlier generated (and potentially
    validated or rejected) patches can help diversify results in a guided fashion.
    We plan to explore how more sophisticated sampling and search (e.g. beam search)
    can improve Passerine's performance.

\end{itemize}







    
    
    



\subsection{Deployment Considerations}

In an eventual deployment, Passerine may be presented
with bugs that are not within reach of its repair capabilities.
Executing Passerine to attempt to repair such bugs increases compute
costs unnecessarily.
To mitigate this issue, we have begun exploring
\emph{repair abstention}, analogous to abstention in 
classification tasks, to enable Passerine to 
abstain from running on bugs that it is unlikely to fix. 
The filtering criteria defined in Section~\ref{sec:challengeset} can be
used as a guide here, but note that, since the ground-truth patch is not
known for unsolved bugs, the same criteria, which depend on properties of the patch, cannot be applied directly. 
Therefore, applying the same filtering criteria requires auxiliary prediction tasks and knowledge about
how similar bugs were solved in the past.



